\documentclass{article}
\usepackage[utf8]{inputenc}
\usepackage[T1]{fontenc}
\usepackage{geometry}
\usepackage{graphicx}
\usepackage{float}
\usepackage{hyperref}
\usepackage[english,indonesian]{babel}

\title{Optimasi Pemrosesan Template LaTeX Menggunakan Pendekatan Berbasis AI}
\author{Academic Support Agent}
\date{\today}

\begin{document}
\maketitle

\begin{abstract}

\end{abstract}

\chapter{PENDAHULUAN}


\chapter{TINJAUAN PUSTAKA}


% Original content (commented out):
% \documentclass[11pt, titlepage]{article}

\usepackage[margin=1in]{geometry}
\usepackage[strict]{changepage}
\usepackage{float}
\usepackage{fancyhdr}
\usepackage{mhchem}
\usepackage{siunitx}
\usepackage{wrapfig, booktabs}
\usepackage{enumitem}
\usepackage{caption}
\usepackage{commath}
\usepackage{amsmath}
\usepackage[hang]{footmisc}
\usepackage{multicol}
\usepackage{amsfonts}
\usepackage{mathrsfs}

\newcommand{\experimentDate}{23 ExampleMonth 2016}
\newcommand{\className}{Class 1A}
\newcommand{\sectionNumber}{1234}
\newcommand{\studentLabNum}{Locker 123}
\newcommand{\experimentNumber}{LAB X}
\author{KERLEGUER Baptiste, LORIDANT Thomas, LOUART Maëlic, PERRIN Vivien}
\newcommand{\authorLastName}{ KERLEGUER, LORIDANT, LOUART, PERRIN}
\title{GENEBIKE}
\newcommand{\experimentShortName}{GENEBIKE}

\date{\parbox{\linewidth}{}}

\pagestyle{fancy}
\fancyhf{}
\rhead{\authorLastName\ \thepage}
\lhead{\experimentShortName}
\cfoot{\className\ -- \experimentNumber}

\usepackage{color}
\usepackage{sectsty}

\definecolor{WordSectionBlue}{RGB}{30, 90, 147}

\allsectionsfont{\color{WordSectionBlue}}

\newcommand{\gpmol}{\si{\gram\per\mol}}

\begin{document}
	\maketitle
    \setcounter{tocdepth}{1}
    \section{Introduction}
    \subsection{Magnetic circuit conception}
    \begin{paragraph}
    Firstly, the magnetic circuit have to be designed. In order to select a geometry for this circuit, a magnetic simulation software was used which is nammed ANSYS Maxwell. These criteria was used to select or not a magnetic circuit:
    \begin{itemize}
    \item Enough important magnetic flow per turn per inductor to reach a 6V no-load voltage with a reasonable number of turn per inductor(less than 80).
    \item Low detent torque.
    \end{itemize}
    To begin, an axial flow geometry was studied but the results of the simulation have showed that the flow criteria can not be reach with this geometry.\\
    Thus a first radial flow geometry was studied but the detent torque was too important and the second criterion was not validated.\\
    Finally a second axial flow geometry using supports in alloy for inductor in order to reduce the detent torque was designed and it was showed that the criteria was validated by simulation.
    \end{paragraph}
    \section{Description}
    
    \section{Error Analysis \& Discussion}
    
    \section{Conclusion}
    
    \newpage
   	\section{Data Tables}
    
    
\end{document}

\end{document}